\subsection{}
\label{I.introduction.2}

El presente volumen constituye el tomo~1 del «\emph{Seminario de Geometría Algebraica del Bois Marie}», cuyos volúmenes siguientes está previsto que aparezcan en la misma serie que este. El objetivo que se propone el \emph{Seminario}, paralelamente al tratado «\emph{Elementos de geometría algebraica}» de J\ptbl\ \textsc{Dieudonné} y A\ptbl\ \textsc{Grothendieck}, es sentar los fundamentos de la geometría algebraica desde los puntos de vista expuestos en esta última obra. La referencia estándar para todos los volúmenes del \emph{Seminario} está constituida por los capítulos~I, II y III de los «\emph{Elementos de geometría algebraica}» (citados como EGA~I, II y III), y se supone que el lector posee el bagaje de álgebra conmutativa y álgebra homológica que estos capítulos requieren\footnote{Véase la Introducción de EGA~I para más precisiones al respecto.}. Además, en cada volumen del \emph{Seminario} se remitirá libremente, en la medida de lo necesario, a volúmenes anteriores del mismo \emph{Seminario} o a otros capítulos publicados o a punto de aparecer de los «Elementos».

Cada \emph{parte} del \emph{Seminario} se centra en un tema principal, indicado en el título del volumen o de los volúmenes correspondientes; el seminario oral abarca generalmente un año académico, a veces más. Las exposiciones comprendidas en cada parte del \emph{Seminario} mantienen generalmente entre sí una estrecha dependencia lógica; en cambio, las diferentes partes del \emph{Seminario} son, en gran medida, lógicamente independientes unas de otras. Así, la parte «Esquemas en grupos» es casi enteramente independiente de las dos partes del \emph{Seminario} que la preceden cronológicamente; en cambio, recurre con frecuencia a los resultados de EGA~IV. He aquí la lista de las partes del \emph{Seminario} que han de aparecer próximamente (citadas en lo sucesivo como SGA~1 a SGA~7):
\begin{enumerateb}
\item[SGA 1.] Recubrimientos étales y grupo fundamental, 1960 y 1961.
\item[SGA 2.] Cohomología local de los haces coherentes y teoremas de Lefschetz locales y globales, 1961/62.
\item[SGA 3.] Esquemas en grupos, 1963 y 1964 (3~volúmenes), en colaboración con M\ptbl\ \textsc{Demazure}.
\item[SGA 4.] Teoría de los topoi y cohomología étale de los esquemas, 1963/64 (3~volúmenes), en colaboración con M\ptbl\ \textsc{Artin} y J.L\ptbl\ \textsc{Verdier}.
\item[SGA 5.] Cohomología $l$-ádica y funciones~$L$, 1964 y 1965 (2~volúmenes).
\item[SGA 6.] Teoría de intersecciones y teorema de Riemann--Roch, 1966/67 (2~volúmenes), en colaboración con P\ptbl\ \textsc{Berthelot} y L\ptbl\ \textsc{Illusie}.
\item[SGA 7.] Grupos de monodromía local en geometría algebraica.
\end{enumerateb}

Tres de estos \emph{Seminarios} parciales fueron dirigidos en \emph{colaboración} con otros matemáticos, que figurarán como coautores en la cubierta de los volúmenes correspondientes. En cuanto a los demás participantes activos del \emph{Seminario}, cuyo papel (tanto desde el punto de vista de la redacción como desde el de la puesta a punto matemática) ha ido creciendo de año en año, el nombre de cada participante figura al comienzo de las exposiciones de las que es responsable como conferenciante o como redactor, y la lista de quienes figuran en un volumen determinado se encuentra
\ifthenelse{\boolean{orig}}
{indicado}
{indicada}
en la portada de dicho volumen.

Conviene dar algunas precisiones sobre las relaciones entre el \emph{Seminario} y los \emph{Elementos}. Estos últimos estaban destinados en principio a ofrecer una exposición de conjunto de las nociones y técnicas consideradas más fundamentales en geometría algebraica, a medida que estas mismas nociones y técnicas se fueran desprendiendo mediante el juego natural de exigencias de coherencia lógica y estética. Desde esta perspectiva, era natural considerar el \emph{Seminario} como una versión preliminar de los \emph{Elementos}, destinada a quedar incorporada casi por entero, tarde o temprano, en estos últimos. Este proceso ya había comenzado en cierta medida hace algunos años, puesto que las exposiciones I a IV del presente volumen SGA~1 están enteramente incorporadas en EGA~IV, y las exposiciones VI a VIII debían estarlo dentro de algunos años en EGA~VI. Sin embargo, a medida que se desarrolla la labor de edificación emprendida en los \emph{Elementos} y en el \emph{Seminario}, y se precisan las proporciones del conjunto, el principio inicial (según el cual el \emph{Seminario} no constituiría más que una versión preliminar y provisional) parece cada vez menos realista, debido, entre otras cosas, a los límites que la previsora naturaleza impone a la duración de la vida humana. Habida cuenta del cuidado prestado generalmente a la redacción de las distintas partes del \emph{Seminario}, sin duda solo habrá motivo para retomar una parte semejante en los \emph{Elementos} (o en los tratados que los sucedan) cuando los progresos posteriores a su redacción permitan introducir en ella mejoras muy sustanciales, al precio de modificaciones bastante profundas. Este es ya el caso del presente seminario SGA~1, como se ha dicho más arriba, y también de SGA~2 (gracias a los resultados recientes de la Sra.~\textsc{Raynaud}). En cambio, nada indica actualmente que vaya a ser así en un futuro próximo para ninguna de las partes SGA~3 a SGA~7 citadas más arriba.

Las referencias dentro del «Seminario de Geometría Algebraica del
\ifthenelse{\boolean{orig}}
{Bois-Marie}
{Bois Marie}» se dan del modo siguiente. Una \emph{referencia interna} a una de las partes SGA~1 a SGA~7 del Seminario se da en la forma III~9.7, donde la cifra III designa el número de la exposición, que figura en la parte superior de cada página de la exposición en cuestión, y~9.7, el número del enunciado (o de la definición, observación, etc.) dentro de la exposición. Cuando sea necesario, pueden utilizarse números decimales más largos, por ejemplo 9.7.1 y 9.7.2, para designar, por ejemplo, las distintas etapas de la demostración de una proposición~9.7. La referencia III~9 designa el parágrafo~9 de la exposición~III. El número de la exposición se omite en las referencias internas a una exposición. Para una \emph{referencia a otra} de las \emph{partes} del \emph{Seminario}, se utilizan las mismas siglas, pero precedidas de la mención de la parte de los SGA de que se trate: SGA~1~III~9.7. Del mismo modo, la referencia EGA~IV~11.5.7 significa: \emph{Elementos de geometría algebraica}, cap\ptbl\ IV, enunciado (o definición
\ifthenelse{\boolean{orig}}{etc...}{etc.})~11.5.7; aquí, la primera cifra arábiga sigue designando el número del parágrafo. Aparte de estas convenciones vigentes en el conjunto de los SGA, la bibliografía relativa a una exposición estará reunida generalmente al final de esta, y dentro de la exposición se remitirá a ella mediante números entre corchetes, como
\ifthenelse{\boolean{orig}}{[3]}{[\textbf{3}]}, según el uso habitual.

Por último, para comodidad del lector, siempre que parezca necesario añadiremos al final de los volúmenes de los SGA un índice de notaciones y un índice terminológico que contenga, cuando proceda, una traducción inglesa de los términos franceses utilizados.

Deseo añadir a esta introducción un comentario extramatemático. En noviembre de 1969 tuve conocimiento de que el Institut des Hautes Études Scientifiques, donde he sido profesor casi desde su fundación, recibía desde hacía tres años subvenciones del Ministerio de las Fuerzas Armadas. Ya como investigador principiante me parecía sumamente lamentable la falta de escrúpulos con que la mayoría de los científicos aceptan colaborar, de una forma u otra, con los aparatos militares. Mis motivaciones de entonces eran esencialmente de naturaleza moral y, por tanto, poco susceptibles de ser tomadas en serio. Hoy adquieren una fuerza y una dimensión nuevas, en vista del peligro de destrucción de la especie humana con que nos amenaza la proliferación de los aparatos militares y de los medios de destrucción masiva de que estos disponen. En otros lugares me he explicado con mayor detalle sobre estas cuestiones, mucho más importantes que el progreso de cualquier ciencia (incluida la matemática); puede consultarse, por ejemplo, a este respecto el artículo de G\ptbl\ Edwards en el n.º~1 de la revista \emph{Survivre} (agosto de 1970), que resume una exposición más detallada de estas cuestiones que había hecho en otro lugar. Así pues, durante tres años me encontré trabajando en una institución mientras esta participaba, sin que yo lo supiera, en una modalidad de financiación que considero inmoral y peligrosa\footnote{Huelga decir que la opinión que acabo de expresar es de mi exclusiva responsabilidad y no de la editorial Springer, que publica el presente volumen.}. Al ser actualmente el único de mis colegas del IHES que mantiene esta opinión, lo que ha condenado al fracaso mis esfuerzos por lograr que se supriman del presupuesto del IHES las subvenciones militares, he tomado la decisión que se imponía: dejo el IHES el 30 de septiembre de 1970 y suspendo asimismo toda colaboración científica con esta institución mientras siga aceptando tales subvenciones. He pedido al Sr\ptbl\ Motchane, director del IHES, que, a partir del 1 de octubre de 1970, el IHES se abstenga de difundir textos matemáticos de los que soy autor o que formen parte del Seminario de Geometría Algebraica del Bois Marie. Como se ha dicho más arriba, la difusión de este seminario será asumida por la editorial Julius Springer, en la serie Lecture Notes. Me complace agradecer aquí a la editorial Springer y al Sr\ptbl\ K\ptbl\ Peters la ayuda eficaz y cortés que me han prestado para hacer posible esta publicación, haciéndose cargo, en particular, de la mecanografía para la reproducción fotoóffset de las nuevas exposiciones añadidas a los antiguos seminarios y de las exposiciones que faltaban en los seminarios incompletos.

Agradezco asimismo al Sr\ptbl\ J.P\ptbl\ Delale, que se encargó de la ingrata labor de compilar el índice de notaciones y el índice terminológico.

\bigskip\hfill Massy, agosto de 1970.

